\subsection{2.5 选择训练数据配比}\label{selecting-a-data-mixture-for-training}

训练期间采样的数据源及其相对比例会显著影响最终模型。在算力预算固定、语料库包含数百个异构数据源的情况下，我们必须确定如何为它们分配权重。

为此，我们首先定义一个目标函数。如图 5 所示，给定一组目标能力（如编码或通用知识），可以得到多个帕累托高效（Pareto-efficient）配比，在其中一个能力无法在不损害另一能力的前提下得到提升。因此，我们必须为这些不同能力分配相对重要性。具体来说，我们选择最小化式 3，即对一系列留出（held-out）NLL 评测的聚合。

优化数据配比需要在感兴趣领域、数据质量、可用的唯一 token 数量，以及一个数据集在收益递减或过拟合之前最多可重复的次数之间进行权衡。

\subsection{2.5.1 优化面临的挑战}\label{challenges-to-optimization}

数据配比的选择对最终模型有重大影响，因此我们希望以有原则的、结构化的方式来解决这一问题。一个有效的技术必须应对以下几项挑战：

• 效用的定义。我们如何具体定义优化所针对的目标？我们的总体目标是开发在真实世界中有效的模型，但对每个候选数据配比都运行一次完整的训练后处理（post-training）流程，代价将高得难以承受。

• 巨大的搜索空间。随着数据集数量增长到数百个，数据配比的空间变得极大，穷举搜索在计算上不可行。

• 与规模相关的影响。某个数据源可能对较小规模的模型特别有用，而另一个数据源在大规模下更有用（Ye et al., 2025）。

• 跨数据集交互。数据集之间在不同方面既互补又可替代。一个数据集的效用无法脱离配比中的其余部分独立衡量。

• 多轮次影响。数据配比无法孤立评估：训练长度很重要（Chen et al., 2026a）。虽然很有吸引力，但对高质量但规模小的数据集赋予过高权重，在训练长度较长时会导致收益递减，因为该数据集最终会被消费多次。在最坏的情况下，这可能导致过拟合。

• 计算成本。训练大量采用不同数据配比的模型很快就会变得昂贵，尤其是在试图探索极高维搜索空间时。

\subsection{2.5.2 基于预测的方法}\label{forecastingbased-approaches}

一类数据选择方法是以显著更小的规模训练大量模型，然后基于这些模型的评测结果构建预测模型。从这些小模型上的拟合函数出发，可以尝试估计全局最优的数据配比，或至少为搜索提供指导。这类方法包括 RegMix（Liu et al., 2025b）等。这些技术的一个关键特征是，通过依赖小模型，它们能够实现快速迭代。

我们尝试了这些方法的变体：基于一个基线配比采样数据配比，并训练了几千个小模型，每个模型的激活参数在 760M 到 4B 之间（按表 1 为 L12-L36 模型）——Web / Code 数据前沿的示例如图 5 所示。对每个模型，我们计算完整的 NLL 评测套件，考察不同数据子集与不同评测之间的相关性——包括聚合评测与单项评测——示例见图 5 和图 7。最终，我们按照 2.3 节中的式 3 为每个配比给出单一得分，并试图对其进行全局优化（最小化）。

实践中尺度依赖的一个实例。在数据配比实验期间，我们发现了一个挑战我们核心缩放假设之一的经验结论：即两个数据配比的相对排序——哪个配比能带来更好的评测表现——会随着算力的增加而保持不变。这一假设通常被非正式地称为排序不变性假设（rank invariance hypothesis）（Liu et al., 2025b），它之所以有吸引力，是因为它表明廉价的的小规模实验可以可靠地指导大规模运行的决策。然而，我们发现情况并非总是如此。

在早前的一代模型中，我们试图衡量编码与 STEM 强化学习对预训练数据配比的敏感度。我们使用不同的配比优化方法和不同的评测权重，产出了两个候选配比：code-heavy-mix（约 50\% 代码，预计在 Coding 类别上表现尤为突出）和 stem-heavy-mix（显著上调 STEM 数据，以针对 STEM 和 Math 类别）。我们先进行了中间规模的验证，随后在每种配比上训练了约 20T tokens 的 23B 激活参数模型（按表 1 大致相当于 L66）。除训练数据源的权重外，其余均未改动。

\pandocbounded{\includegraphics[keepaspectratio,alt={image}]{images/d344b0ec115e0d9f8ceacd339f966fd02475ebc4532ad2d0459c8be33001994a.jpg}}

图 6. 数据配比缩放中的排序非不变性。（左）：在较小规模下，stem-heavy-mix 在留出 STEM NLL 评测上表现更好，但在更大（23B）规模下 code-heavy-mix 表现更好。（右）：两种 23B 激活参数模型在不同数据配比下整个训练过程中 STEM 评测的 NLL。

\pandocbounded{\includegraphics[keepaspectratio,alt={image}]{images/58df6ac093c1bd77ffc7af0cfaad2b6d80c74ce918f6a591bb870f73f05b61c9.jpg}}

结果如图 6 所示。与小规模结果一致，训练早期 stem-heavy-mix 在 NLL STEM 评测上优于 code-heavy-mix。然而，STEM 评测曲线在训练中途发生交叉，code-heavy-mix 最终反超 stem-heavy-mix。这与"数据配比的相对排名会随训练算力扩大而保持不变"的假设相矛盾。

通过检查各数据源的训练与验证损失趋势，我们识别出两个数据源：它们包含高质量的 STEM 内容，但比其他数据集有更多模糊重复，且内容多样性较低。这两个数据源在 stem-heavy-mix 中的权重为 11.8\%，而在 code-heavy-mix 中仅为 0.3\%。我们的假设是：这些数据源对较小模型非常有用，但其多样性不足在更大规模下本质上帮助有限。

基于这一结果，除了候选数据配比在固定实验规模下的表现，我们还更加重视其缩放性能。此外，在后续的数据配比研究中，我们通过缩放阶梯方法极其仔细地关注了数据配比的缩放特性。

\subsection{2.5.3 最终配比选择}\label{final-mix-selection}

由于数据集数量众多，且我们希望在不同规模下对其进行验证，我们采用了一种分层方法来改进此前的数据配比。我们将数据划分为约 10 个类别，如代码、STEM、PDF 和通用网页。然后交替进行局部与全局优化：

• 局部搜索。我们保持所有高层类别固定，在单个子集内局部调整权重——例如，在代码文件、PR 和 commit 之间调整权重，同时保持通用网页、STEM 等权重不变。

• 全局搜索。我们保持每个高层分箱的相对构成固定，调整不同子集之间的相对权重（例如，上调 STEM 和 PDF，降低所有其他组件的权重）。图 7 展示了这一过程。

研究生水平物理 NLL

图 7. 对于一套由第三方创建的研究生水平物理评测集，我们展示了不同数据集的相对贡献（NLL 越低越好）。我们采样不同的数据配比，分别训练并评估所得模型，然后考察最终 NLL 与配比数据构成之间的相关性。我们的 PDF 数据集、从网页中提取的数学与 STEM 数据呈现出明显的正向趋势。通用网页语料基本为中性，而在配比中增加更多代码（而非其他数据）似乎对物理并无帮助。

\pandocbounded{\includegraphics[keepaspectratio,alt={image}]{images/933b3af442479d3aaf99faac29862d775576881234f8a12092bf9e8054367864.jpg}}

在整个搜索过程中，我们将任一数据集的重复次数上限设为 8，以防止过拟合和反复循环同一数据带来的收益递减（Chen et al., 2026a）。

在进行了多轮全局与局部数据优化的迭代之后，我们选出几个最小化 NLL 目标函数的候选配比。我们以约 2.8$\times$ 倍于全局混配所用算力的规模训练了这些模型，并选出最佳候选。在对数据配比流程进行这些调整后，我们在这一放大验证步骤中看到了良好的缩放表现：最优候选不再随规模变化。

表 5 总结了最终构成。每个源家族包含多个数据集，各数据集的轮数各不相同。某些数据集可能同时属于多个类别（例如一个法语的 STEM 数据集），在表中我们按自上而下的顺序将每个数据集归入第一个匹配的行。在最终配比中，我们得到了约 3000 亿高质量数学 tokens，其平均采样倍数为 5.28$\times$，是所有源家族中最高的。不过，训练配比的主体是代码数据，达到 16.4T 训练 tokens，平均略多于 2 轮。网页文本和 PDF 数据平均各被看到不到一次（0.55$\times$ 和 0.53$\times$），这意味着即使在整个 30T-token 的训练中，这些数据源的全部可用语料也未被用完。需要说明的是，对于 PDF 和网页文本，其中的数学与 STEM 部分被计入相应的学科类别。我们的多语言语料是被降采样最激进的：可用 8.1T 唯一 tokens，但最终配比仅消耗 0.5T（0.06$\times$）——同样需要说明，特定领域的多语言数据已被包含在不同的类别中。

表 5. MAI-Base-1 预训练数据构成。唯一 tokens（Unique tokens）为各源家族去重后的 token 数；训练 tokens（Training tokens）为整个训练过程中从该家族消耗的 token 数；平均轮数（Avg. Epochs）为训练 tokens 与唯一 tokens 之比，大于 1 表示重复采样。

{\def\LTcaptype{none} % do not increment counter
\begin{longtable}[]{@{}|l|l|l|l|l|@{}}
\toprule\noalign{}
\endhead
\bottomrule\noalign{}
\endlastfoot
\hline
数据源家族 & 唯一 tokens (T) & 训练 tokens (T) & 配比百分比
(\%) & 平均轮数 \\
\hline
代码 & 7.4 & 16.4 & 54.6 & 2.22$\times$ \\
\hline
STEM & 2.2 & 4.7 & 15.8 & 2.17$\times$ \\
\hline
数学 & 0.3 & 1.6 & 5.4 & 5.28$\times$ \\
\hline
书籍与期刊 & 0.6 & 0.9 & 3.1 & 1.65$\times$ \\
\hline
PDF & 2.7 & 1.4 & 4.7 & 0.53$\times$ \\
\hline
网页文本 & 8.1 & 4.5 & 14.9 & 0.55$\times$ \\
\hline
多语言（其他） & 8.1 & 0.5 & 1.6 & 0.06$\times$ \\
\hline
总计 & 29.2 & 30.0 & 100.0 & 1.03$\times$ \\
\hline
\end{longtable}
}

我们仍在持续迭代数据配比目标、构成该目标的基于 NLL 的评测，以及用于提出新候选配比的算法，因为每个组件都对最终数据配比有重大影响。理解端到端的缩放行为——以及找出在预训练期间最可靠地预测下游性能的指标——仍是一个基础性的开放挑战，我们认为其中仍有很大的提升空间。

\subsection{2.5.4 训练中期的数据配比}\label{midtraining-data-mixture}

预训练之后，MAI-Base-1 按 2.6.1 节的描述分两个连续阶段进行训练中期（mid-training）训练。训练中期数据完全来自预训练语料（2.4 节），不引入任何新的或合成的数据源。我们通过对现有配比进行过滤和重新加权来整理出质量更高的子集，然后以更长的序列长度重新打包结果，用于上下文扩展。

总体而言，我们进一步将配比向 STEM、数学和代码倾斜，为推理 RL 的爬升奠定坚实基础（Wang et al., 2025b）。在预训练配比的基础上，我们将 STEM/数学数据的比例提高到 35\%，代码保持在 55\%，其余 10\% 分配给背景数据源。在每个类别内部，我们根据单源微退火（microanneals）（Team OLMo et al., 2024）局部调整权重，优化与预训练配比搜索相同的基于 NLL 的目标，并针对上下文扩展加入长上下文 NLL 任务（附录 C.2）。此外，我们还对部分选定的 STEM 和代码数据源应用了质量过滤，具体如下。

STEM 推理数据。为提升推理任务的表现，我们主要使用（Essential AI et al., 2025）提出的内容质量分类法对预训练 PDF 语料中的文档进行过滤。我们定义了一个 Bloom Analyze 启发式规则：保留技术正确性高、推理深度至少为中等、且 Bloom 认知加工水平达到或高于 Analyze（分析）层的文档（Anderson and Krathwohl, 2001）。同时去除存在提取痕迹或缺少主要内容的文档。该过滤器保留具有结构化推理和分析的文档，而非简单的事实陈述。

代码数据。训练中期的代码数据还要经过两个额外步骤。首先，我们按质量分档（quality-bin）进行文件扩展名过滤：仓库根据仓库级质量指标分为三个层级。例如，常见的网页资源（HTML、CSS、SVG）在顶级仓库中会保留，因为它们通常属于较大的前端项目；但在较低层级仓库中会被移除，因为这些仓库中它们主要是低质量的独立页面。其次，我们在预训练已有的仓库级格式之外引入文件级文档格式，将仓库理解和单个文件理解视为互补任务，这提升了内部文件局部 NLL 基准上的泛化表现。

记忆感知的轮次上限。除了质量过滤，我们还在训练中期限制数据源的暴露量，以避免对预训练结束时已被高度记忆的数据源过拟合。我们通过每个数据源的验证代理指标来估计这一点：两个预训练 checkpoint 之间验证损失的改善中，来自以接近确定性预测出的 tokens 的比例
\(\left( \mathrm { N L L } < 0 . 0 1 \right)\) 。该比例越高，说明 NLL 的显著降低来自记忆化或高度重复的结构。在训练中期，我们对这类数据源设定更严格的轮次上限，同时允许数值较小的数据源获得相应更多的暴露。

长上下文数据配比。训练中期分两个阶段将 MAI-Base-1 的上下文窗口扩展到 64k 再到 256k，详见附录 C。两个阶段都在更长的序列长度下重新打包数据，而不修改配比权重。这样做可以最大程度地减少分布偏移，同时显著减少现有高质量长上下文文档（如教学或学术文本）被截断的情况。

\subsection{2.6 训练方案}\label{training-recipe}

MAI-Base-1 分三个阶段训练，序列长度逐步增加：预训练、训练中期阶段 1 和训练中期阶段 2。下面介绍各阶段的训练方案。

\subsection{2.6.1 训练阶段}\label{training-phases}

表 6 列出了从 2.5.3 节所述数据配比中采样的 token 数、上下文长度以及各阶段使用的 GPU 数量。预训练和训练中期 1 使用专家并行度 64 搭配 ZeRO-2，训练中期阶段 2 启用 ZeRO-3（FSDP）。训练中期阶段使用上下文并行。训练使用我们自研的大规模训练框架（2.8 节）。

表 6. MAI-Base-1 各阶段的训练规格。

{\def\LTcaptype{none} % do not increment counter
\begin{longtable}[]{@{}|l|l|l|l|@{}}
\toprule\noalign{}
\endhead
\bottomrule\noalign{}
\endlastfoot
\hline
阶段 & Tokens & 上下文长度 & GB200 GPU 数 \\
\hline
预训练 & 30 T & 16,384 & 8,192 \\
\hline
训练中期 1 & 3.4T & 65,536 & 8,192 \\
\hline
训练中期 2 & 150 B & 262,144 & 4,096 \\
\hline
\end{longtable}
}

\subsection{2.6.2 超参数}\label{hyperparameters}

MAI-Base-1 使用 AdamW 优化（Loshchilov and Hutter, 2017）
\(( \beta _ { 1 } = 0 . 9 5 , \beta _ { 2 } = 0 . 9 2 5 , \varepsilon = 1 0 ^ { - 8 } )\)
，权重衰减恒为 0.1，其中注意力权重（0.01）和嵌入权重（0.005）使用较小的权重衰减，以避免对受益较少的参数过度正则化。全局梯度范数裁剪到 1.0。学习率经过约 12B tokens 的线性预热后，按余弦调度从峰值 \(2 \times 1 0 ^ { - 4 }\) 衰减到最小值 2 $\times$ 10−5——最终与峰值的比值为 0.1$\times$，而不是更常见的 0.01$\times$，因为我们发现学习率衰减更少可以改善 RL 之后的结果。我们在每层输出上、残差相加之前应用 0.15 的 dropout。虽然据我们所知高 dropout 并非标准做法，但我们发现它与权重衰减一起提供了互补的正则化效果，并且观察到它能提升缩放阶梯评测任务上的表现。默认情况下，模型权重从均值为零、标准差为 0.02 的正态分布初始化，残差相加之前的输出投影权重按模型中残差连接总数的平方根的倒数进行缩放（Radford et al., 2019）。

图 8. (a) 随机初始化时，注意力层最初产生坍缩的表示——本质上是在序列维度上执行因果平均池化。注意力贡献中不同行之间存在高相关性，这会降低 token 表示的多样性，并导致后续 MoE 层高度不平衡。(b) 将注意力层的输出初始化为零后，MoE 层中的 token 分布更均衡。专家的违反率（violation rate）是其与各专家平均 token 数的偏差比例。

\pandocbounded{\includegraphics[keepaspectratio,alt={image}]{images/8703e55c198595bd088d4b3a717378b29c1844073cef9b56b5e26472f5a526e4.jpg}}

所有阶段使用 134M tokens 的全局 batch size。训练中期阶段继承预训练的超参数设置，除非另有说明。训练中期阶段 1 使用余弦衰减学习率调度，峰值学习率为 \(2 \times 1 0 ^ { - 5 }\)，最低为 \(1 \times 1 0 ^ { - 5 }\)，无预热。训练中期阶段 2 使用恒定学习率 \(1 \times 1 0 ^ { - 6 }\)，同样无预热。

注意力初始化。初始化时，注意力 softmax 接近均匀分布，实际上等同于在因果约束下对 tokens 进行平均池化，如图 8 所示。这种行为降低了 token 表示的多样性，并可能导致后续 MoE 层出现高度不平衡的路由，且随深度增加而加剧。为解决这一问题，我们将注意力输出初始化为零，具体做法是将输出 RMSNorm 的增益（gains）设为零。这意味着模型初始表现为一堆应用于单个 token 的前馈层，而注意力层捕获的跨 token 交互会在训练过程中逐渐生效。我们确认，该方法显著降低初始峰值不平衡，从而带来了 EG 测量的质量提升（见 2.2.2 节）。

\subsection{2.6.3 数值精度}\label{numerical-precision}

模型训练中，权重和激活的默认数据类型为 BF16。前向传播中的 GEMM 使用 FP8 E4M3，数据梯度使用 FP8 E5M2，权重梯度的计算使用 BF16，并采用 FP32 梯度累积（Micikevicius et al., 2017, 2022）。所有 FP8 运算使用延迟缩放，历史窗口为 1024 步的绝对值最大值。

我们在许多对数值敏感的位置使用 FP32。在激活方面，包括所有 softmax 之前的激活（注意力分数、MoE router logits 和最终输出 logits）、MoE combine，以及从输入嵌入到输出的整个残差流。在参数方面，包括嵌入权重、RMSNorm 权重和 router 权重。优化器中的主参数和动量缓冲区（buffers）保持 FP32，所有 AdamW 计算也是如此。数据并行 all-reduce 和微批次梯度累积缓冲区以 FP32 运行。MoE 路由和最终词表权重 GEMM 也均保持 FP32。

图 9. MAI-Base-1（35B 激活 / 1T 总参数）在 30T tokens 上的预训练损失曲线，未做平滑处理。损失在几次早期尖峰后迅速恢复。整个训练过程中没有跳过任何训练批次，也没有对训练配置进行任何手动干预。

\pandocbounded{\includegraphics[keepaspectratio,alt={image}]{images/310d9d10402d320ea96715cba3596f7755b52d28d42ec2e845a367cf8d2ede2d.jpg}}

我们尽可能将类型转换与相邻运算（如 RMSNorm）融合。当梯度从高精度流向低精度时，应用随机舍入（stochastic rounding）（Zamirai et al., 2020; Ozkara et al., 2025）。这是因为层内计算精度低于残差流的精度，因此反向传播中需要对梯度进行降精度转换。

\subsection{2.6.4 训练损失曲线}\label{training-loss-curve}

MAI-Base-1 的训练损失曲线见图 9，未做平滑处理。训练早期出现了几次尖峰，但损失值每次都迅速恢复，我们没有干预便继续了训练。通过运行时日志，我们发现这些尖峰主要影响代码数据集，并与无丢弃（dropless）路由下的高专家不平衡相关：即 tokens 过度集中在少数几个专家中。
